\input{../tex-inputs/latex-leseansicht-vorspann}

\section[Rainer Maria Rilke: Widmungsexemplar von Lieder am Volke geschenkt, 1. 4. 1896]{L04539 Rainer Maria Rilke: Widmungsexemplar von Lieder am Volke geschenkt, 1. 4. 1896}
\nopagebreak\mylabel{L04539v}
\rehead{ }\normalsize\beginnumbering\briefempfaengerindex{Schnitzler, Arthur@\textsc{Schnitzler, Arthur}!zzzRilke, Rainer Maria@\emph{von Rainer Maria Rilke}!1896-04-013@{1. 4. 1896}|(be}
\toendnotes[C]{\smallbreak\pagebreak[2]}
\correspDesc{Versand  durch Rainer Maria Rilke am 1. 4. 1896 in Prag
\newline{}Erhalt  durch Arthur Schnitzler im Zeitraum [2. 4. 1896 – 6. 4. 1896?] in Wien}\toendnotes[C]{\smallbreak}
\buchAlsQuelle{\emph{Rainer Maria Rilke und Arthur Schnitzler. Ihr Briefwechsel.}. Mit Anmerkungen herausgegeben von Heinrich Schnitzler In: \emph{Wort und Wahrheit}, Jg. 13, Nr. 4, April 1958, S. 282–298, hier S. 293.}\toendnotes[C]{\smallbreak}
\pstart
           \noindent{}\raggedleft{}{\pb}{[}1. April
                  1896{]}\pend
           
\pstart
           Meister Arthur Schnitzler mit dem
                  Ausdrucke der größten aufrichtigsten Verehrung.\pwindex{Wegwarten@\emph{Wegwarten}|pwv}\pend
           \pstart \spacefill\mbox{René Maria Rilke.}\pend{}\selectlanguage{ngerman}\endnumbering\briefempfaengerindex{Schnitzler, Arthur@\textsc{Schnitzler, Arthur}!zzzRilke, Rainer Maria@\emph{von Rainer Maria Rilke}!1896-04-013@{1. 4. 1896}|)be}\mylabel{L04539h}
\begin{anhang}
\end{anhang}\newcommand{\dateiname}{L04539}\newcommand{\titel}{Rainer Maria Rilke: Widmungsexemplar von Lieder am Volke geschenkt, 1. 4. 1896}\newcommand{\editorInnen}{Herausgegeben von Jahnke, SelmaMüller, Martin Anton}\input{../tex-inputs/latex-leseansicht-abspann}
